O \gls{mev}, desenvolvido pela \emph{Northrop Grumman}, é uma espaçonave voltada à extensão da vida útil de satélites em órbita geossíncrona.
Lançado em 9 de outubro de 2019, o \emph{MEV-1} realizou, em fevereiro de 2020, sua primeira operação de acoplamento com um satélite-alvo (\emph{Intelsat 901}), com a finalidade de prolongar sua permanência operacional \cite{nytMEV1}.
Em abril de 2025, a missão foi concluída após o reposicionamento do \emph{Intelsat 901} para uma órbita cemitério \cite{intelsat901}.

A extensão de missão foi viabilizada por um mecanismo de captura acoplado à interface do motor de apogeu do satélite-alvo, permitindo ao \gls{mev} assumir funções de controle de atitude e manutenção orbital.
Esse tipo de serviço reduz a necessidade de substituição imediata por novos satélites e pode contribuir para estratégias de gestão de ativos espaciais.
Além disso, a missão constitui uma demonstração de capacidades de operações de proximidade e acoplamento em órbita, relevantes para o campo de \gls{oos} \cite{mev1}.
